% ============================================================================
% APÊNDICE A — Instalação do ambiente
% ============================================================================
\chapter{Instalação e Configuração do Ambiente}
\label{apendice:instalacao}

Este apêndice configura o ambiente completo que você usará em todos os
capítulos. Siga na ordem — cada passo verifica o anterior.

\section{Editor de código: VS Code}

\begin{enumerate}[leftmargin=*, itemsep=0.3em]
  \item Baixe e instale o VS Code: \url{https://code.visualstudio.com};
  \item Extensões recomendadas: \emph{ESLint}, \emph{Prettier},
  \emph{TypeScript (nativo)}, \emph{Tailwind CSS IntelliSense},
  \emph{Prisma}, \emph{Docker}, \emph{GitLens} e o tema de sua preferência;
  \item Ative \emph{Format on Save} (Ctrl+Shift+P $\rightarrow$ Settings
  $\rightarrow$ \texttt{editor.formatOnSave}).
\end{enumerate}

\section{Node.js e npm}

\begin{enumerate}[leftmargin=*, itemsep=0.3em]
  \item Baixe o instalador LTS em \url{https://nodejs.org} (recomendamos a
  versão 24 LTS, usada neste livro);
  \item Verifique no terminal:
\end{enumerate}

\begin{lstlisting}[style=livro, language=Bash, caption={Verificação do Node.}, label={list:inst-node}]
node --version   # v24.x ou superior
npm --version    # 11.x
\end{lstlisting}

\begin{caixadica}
No Windows, o terminal recomendado é o \textbf{Git Bash}, instalado junto com
o Git (abaixo). Ele permite os comandos \texttt{ls}, \texttt{curl} e scripts
como os deste livro sem configuração extra.
\end{caixadica}

\section{Git e GitHub}

\begin{enumerate}[leftmargin=*, itemsep=0.3em]
  \item Instale o Git: \url{https://git-scm.com};
  \item Crie uma conta em \url{https://github.com};
  \item Configure sua identidade (uma vez por máquina) e um token de acesso
  pessoal (Personal Access Token) com escopo \texttt{repo}.
\end{enumerate}

\begin{lstlisting}[style=livro, language=Bash, caption={Configuração inicial do Git.}, label={list:inst-git}]
git config --global user.name "Seu Nome"
git config --global user.email "voce@exemplo.com"
git --version
\end{lstlisting}

\section{Docker Desktop}

\begin{enumerate}[leftmargin=*, itemsep=0.3em]
  \item Instale o Docker Desktop: \url{https://www.docker.com/products/docker-desktop/};
  \item Verifique:
\end{enumerate}

\begin{lstlisting}[style=livro, language=Bash, caption={Verificação do Docker.}, label={list:inst-docker}]
docker --version
docker compose version
\end{lstlisting}

\section{PostgreSQL (opções)}

\begin{itemize}[leftmargin=*, itemsep=0.3em]
  \item \textbf{Opção A (recomendada)}: rode via Docker sempre que precisar
  (capítulo~\ref{cap:docker}); 
  \item \textbf{Opção B}: instale o PostgreSQL 18 local
  (\url{https://www.postgresql.org/download/}) com \texttt{psql} e \texttt{pgAdmin}.
\end{itemize}

\begin{lstlisting}[style=livro, language=Bash, caption={Subindo um PostgreSQL de teste.}, label={list:inst-pg}]
docker run --name pg-livro -e POSTGRES_USER=app -e POSTGRES_PASSWORD=senha \
  -e POSTGRES_DB=app -p 5432:5432 -d postgres:18
\end{lstlisting}

\section{Ferramentas de teste de API}

\begin{itemize}[leftmargin=*, itemsep=0.3em]
  \item \textbf{Insomnia} (\url{https://insomnia.rest}) ou \textbf{Bruno}
  (\url{https://www.usebruno.com}, open-source e recomendado);
  \item O \texttt{curl} do Git Bash também basta para testes rápidos.
\end{itemize}

\section{Verificação final: o teste do semáforo}

Rode tudo de uma vez:

\begin{lstlisting}[style=livro, language=Bash, caption={Checklist do ambiente.}, label={list:inst-check}]
node --version && npm --version && git --version && docker --version
echo "Ambiente pronto!"
\end{lstlisting}

\begin{caixadica}
Cada capítulo informa quais ferramentas usa. Você não precisa instalar nada
além deste apêndice agora — o restante (Redis, Playwright, etc.) vem com o
próprio projeto quando necessário.
\end{caixadica}
